\section{Evaluation Plan and Implementation Roadmap}

\subsection{Evaluation Dimensions}

\begin{longtable}{p{0.25\textwidth}p{0.45\textwidth}p{0.20\textwidth}}
\toprule
\textbf{Dimension} & \textbf{Metric} & \textbf{MVP target} \\
\midrule
Conversion correctness & Exact structural equivalence after conversion; round-trip consistency where possible. & \metric{>= 95\%} on curated tests. \\
Validation accuracy & Correct detection of malformed rows, missing fields, type mismatches, invalid schemas. & \metric{>= 90\%}. \\
RAG relevance & Reviewer judgment of retrieved schema/docs/examples in top-k results. & \metric{>= 80\%} relevant in top 5. \\
Retrieval ranking & Recall@k, nDCG@10, MRR, and schema-hit rate on a curated query set with expected source chunks. & \metric{Beat hybrid baseline before enabling advanced modules}. \\
Graph retrieval utility & Correct entity linking, schema-field match accuracy, and multi-hop answer support from graph evidence. & \metric{>= 85\%} correct links on demo schema set. \\
Graph-NER quality & Entity-level precision/recall/F1, concept-normalization accuracy, relation extraction F1, PHI recall, and reviewer correction rate. & \metric{High recall for PHI; >= 85\% F1 for demo entity types}. \\
Representation learning & Compare pretrained embeddings against SSL-adapted embeddings using Recall@k, nDCG@10, cluster purity, schema-hit rate, and nearest-neighbor review. & \metric{Promote only with statistically meaningful gain}. \\
OCR/document AI & Character error rate, word error rate, key-value extraction F1, table-cell F1, checkbox accuracy, bounding-box provenance, and human correction rate. & \metric{Report per document type and language; review low-confidence fields}. \\
Segmentation & Dice, IoU, HD95, boundary F1, uncertainty quality, reviewer correction time, and modality/device subgroup checks. & \metric{Research demo only; compare against baseline}. \\
Object detection/tracking & Detection mAP, sensitivity, FROC, false positives per image, tracking IDF1/HOTA, temporal consistency, and latency. & \metric{Research demo or design report; no diagnostic claim}. \\
Japan-market readiness & APPI/MHLW/PMDA checklist completion, JP Core/FHIR mapping coverage, Japanese terminology alias coverage, and audit/approval evidence completeness. & \metric{Checklist complete for pilot proposal}. \\
Explanation quality & Human rating for clarity, usefulness, and faithfulness to retrieved context and validation report. & \metric{Average >= 4/5}. \\
Workflow success & Tasks completed without incorrect routing or stalled state. & \metric{>= 85\%} for MVP scenarios. \\
Human review effectiveness & Risky or ambiguous transformations correctly trigger review. & \metric{No critical risky action without review}. \\
Latency & End-to-end response time for simple and complex workflows. & \metric{< 3 sec simple; < 15 sec demo workflow}. \\
Security & Prompt injection and unsafe tool-call tests blocked or flagged. & \metric{No unapproved destructive action}. \\
\bottomrule
\end{longtable}

\subsection{Test Dataset Plan}

\begin{itemize}
  \item Valid JSON, YAML, and CSV examples for basic conversion.
  \item Nested JSON objects and arrays for JSON-to-YAML and YAML-to-JSON conversion.
  \item CSV files with missing values, malformed rows, quoted commas, numeric fields, date fields, and inconsistent columns.
  \item Schema cases with required fields, optional fields, enumerations, type mismatches, and duplicate keys.
  \item Prompt-injection strings embedded inside CSV cells and documentation chunks.
  \item Synthetic/Japanese-style clinical forms, referral letters, lab reports, consent forms, and checkbox/table documents for OCR.
  \item Public or synthetic DICOM samples with metadata, reports, segmentation masks, and source image/frame IDs.
  \item Public medical imaging examples for segmentation, detection, and short tracking/video demonstrations when licensing permits.
  \item Japan-market fixtures: JP Core/FHIR examples, Japanese terminology aliases, APPI-sensitive-field cases, and approval/audit workflows.
\end{itemize}

\subsection{Roadmap Summary}

\begin{figure}[H]
  \centering
  \begin{tikzpicture}[
  font=\scriptsize,
  >=Latex,
  week/.style={draw=sfdark!35, fill=sfgray, minimum width=0.78cm, minimum height=0.42cm, align=center, inner sep=1.5pt},
  phase/.style={draw, rounded corners=4pt, minimum height=0.55cm, align=center, inner sep=3pt},
  line/.style={line width=0.8pt}
]

\foreach \w in {1,...,12} {
  \node[week] (w\w) at ({(\w-1)*0.82},0) {W\w};
}

\node[phase, fill=sfblue!8, draw=sfblue, minimum width=0.78cm] at (0,-0.82) {Scope\\Design};
\node[phase, fill=sfcyan!8, draw=sfcyan, minimum width=2.35cm] at (1.64,-1.38) {Core Data\\Workflow};
\node[phase, fill=sfgreen!8, draw=sfgreen, minimum width=1.55cm] at (4.10,-1.94) {Graph-NER\\Retrieval};
\node[phase, fill=sfpurple!8, draw=sfpurple, minimum width=1.55cm] at (5.74,-2.50) {SSL +\\GraphRAG};
\node[phase, fill=sforange!12, draw=sforange, minimum width=1.55cm] at (7.38,-3.06) {GCP\\MLOps};
\node[phase, fill=sfred!8, draw=sfred, minimum width=1.55cm] at (9.02,-3.62) {Evaluate\\Submit};

\draw[line, sfblue] (0,-0.22) -- (0,-0.55);
\draw[line, sfcyan] (0.82,-0.22) -- (0.82,-0.55) -- (2.46,-0.55) -- (2.46,-0.22);
\draw[line, sfgreen] (3.28,-0.22) -- (3.28,-1.11) -- (4.10,-1.11) -- (4.10,-0.22);
\draw[line, sfpurple] (4.92,-0.22) -- (4.92,-1.67) -- (5.74,-1.67) -- (5.74,-0.22);
\draw[line, sforange] (6.56,-0.22) -- (6.56,-2.23) -- (7.38,-2.23) -- (7.38,-0.22);
\draw[line, sfred] (8.20,-0.22) -- (8.20,-2.79) -- (9.02,-2.79) -- (9.02,-0.22);

\node[align=left, text width=9.8cm, anchor=west] at (0,-4.45) {\textbf{Final milestone:} governed healthcare AI workflow prototype with rule-based validation, Graph-NER, advanced retrieval, self-supervised embedding evidence, explainable outputs, GitHub Actions + GCP MLOps path, monitoring, audit, and executive demo.};

\end{tikzpicture}

  \caption{Twelve-week roadmap summary for the full OJTFlow prototype.}
  \label{fig:roadmap}
\end{figure}

\begin{longtable}{p{0.23\textwidth}p{0.15\textwidth}p{0.34\textwidth}p{0.16\textwidth}}
\toprule
\textbf{Phase} & \textbf{Duration} & \textbf{Main tasks} & \textbf{Deliverables} \\
\midrule
Phase 0: Scope and architecture & Week 1 & Finalize medical scope, Japan-market boundary, governance boundary, datasets, module contracts, risk register, and evaluation criteria. & Approved backlog and architecture baseline. \\
Phase 1: Core data workflow & Weeks 2--4 & Build parser, profiler, rule-based conversion, schema validation, MCP-compatible tool contracts, workflow state, and human review triggers. & Stable core MVP with traceable workflow. \\
Phase 2: Graph-NER and retrieval foundation & Weeks 5--6 & Build medical entity extraction, normalization, graph storage, hybrid search, vector collection, HyDE experiment, and retrieval benchmark. & Graph-NER prototype and measured retrieval baseline. \\
Phase 3: Research acceleration & Weeks 7--8 & Add self-supervised embedding adaptation, reranker interface, GraphRAG context, evidence packaging, supported-claim mapping, multimodal evidence schema, and explanation metrics. & SSL, retrieval, GraphRAG, and evidence report. \\
Phase 4: Japan multimodal and GCP MLOps & Weeks 9--10 & Add OCR/document AI, DICOM metadata flow, SAM/MedSAM or nnU-Net segmentation proof of concept, JP Core/Japan governance checklist, GitHub Actions + GCP deployment blueprint, and model/index versioning. & Multimodal prototype and MLOps package. \\
Phase 5: Monitoring, evaluation, and submission & Weeks 11--12 & Run integrated evaluation, multimodal quality checks, security/privacy hardening, monitoring dry run, documentation, final demo, proposal package, slides, and post-MVP roadmap. & Submission-ready research and industrial prototype. \\
\bottomrule
\end{longtable}

\subsection{MVP Backlog}

\begin{itemize}
  \item Backend conversion API for JSON, YAML, and CSV.
  \item Pydantic-based validation and schema inference.
  \item Agent workflow state and orchestrator.
  \item RAG knowledge base with schemas and transformation examples.
  \item Retrieval benchmark with hybrid search, reranking experiment, and source-attribution checks.
  \item Optional HyDE and schema-graph prototype if baseline retrieval leaves measurable gaps.
  \item Optional self-supervised embedding adaptation experiment using schema/document/conversion pairs.
  \item Graph-NER prototype for medical entity extraction, normalization, relation extraction, and schema linking.
  \item OCR/document AI prototype for scanned clinical-style forms with page/box provenance and human correction.
  \item DICOM metadata and visual-evidence object design for image/video workflows.
  \item SAM/MedSAM or nnU-Net segmentation proof of concept on public medical images.
  \item Object detection/tracking experiment design with FROC/mAP/IDF1-style evaluation plan.
  \item JP Core/Japan FHIR mapping notes, Japanese terminology aliases, and APPI/MHLW/PMDA governance checklist.
  \item MCP server for conversion and validation tools.
  \item Human review endpoint for ambiguous cleaning plans.
  \item Simple UI with input, output, explanation, and agent progress timeline.
  \item Docker Compose deployment.
  \item Test dataset suite and evaluation report.
\end{itemize}
